\chapter{Testing Specification}\label{ch:testing}\label{sim:ch:testing}

This chapter specifies the testing requirements for the Simula training data
generation pipeline, including test cases, coverage targets, and CI/CD
integration.

\section{Test coverage requirements}

\begin{table}[htbp]
\centering
\caption{Simula test coverage targets.}
\label{tab:simula-coverage}
\footnotesize
\begin{tabularx}{\linewidth}{@{}l r r X@{}}
\toprule
\textbf{Component} & \textbf{Line \%} & \textbf{Branch \%} & \textbf{Notes} \\ \midrule
Taxonomy engine & 90 & 85 & Core logic \\
Complexity scorer & 85 & 75 & \\
Critic evaluator & 85 & 75 & \\
Diversity calculator & 85 & 75 & \\
Adaptive features & 80 & 70 & Fallback paths critical \\
\textbf{Overall minimum} & \textbf{85} & \textbf{75} & Enforced in CI \\
\bottomrule
\end{tabularx}
\end{table}

\section{Test fixtures}

Fixtures are located in \texttt{tests/fixtures/simula/}.

\subsection{Taxonomy fixtures}

\begin{itemize}
  \item Complete taxonomy tree with 25+ nodes
  \item Shallow taxonomy (depth 2) for edge cases
  \item Deep taxonomy (depth 6) for performance testing
  \item Malformed taxonomy for error handling
\end{itemize}

\subsection{Training example fixtures}

\begin{table}[htbp]
\centering
\caption{Training example fixture requirements.}
\footnotesize
\begin{tabularx}{\linewidth}{@{}l l X@{}}
\toprule
\textbf{Category} & \textbf{Count} & \textbf{Coverage} \\ \midrule
Easy examples & 20 & Single-table, simple predicates \\
Medium examples & 20 & Joins, grouping, ordering \\
Hard examples & 10 & CTEs, window functions, subqueries \\
Invalid examples & 10 & Should be rejected by critic \\
HANA-specific & 10 & Spatial, temporal, full-text types \\
\bottomrule
\end{tabularx}
\end{table}

\section{Unit test cases}

\subsection{Taxonomy traversal}

\begin{tabularx}{\linewidth}{@{}l l X@{}}
\toprule
\textbf{ID} & \textbf{Priority} & \textbf{Description} \\ \midrule
SIM-TAX-001 & P1 & Traverse taxonomy depth-first \\
SIM-TAX-002 & P1 & Calculate node weights correctly \\
SIM-TAX-003 & P1 & Select node based on complexity target \\
SIM-TAX-004 & P2 & Handle empty taxonomy gracefully \\
SIM-TAX-005 & P2 & Handle malformed taxonomy \\
\bottomrule
\end{tabularx}

\subsection{Complexity scoring}

\begin{tabularx}{\linewidth}{@{}l l X@{}}
\toprule
\textbf{ID} & \textbf{Priority} & \textbf{Description} \\ \midrule
SIM-CMP-001 & P1 & Calculate Elo rating update \\
SIM-CMP-002 & P1 & Classify EASY/MEDIUM/HARD correctly \\
SIM-CMP-003 & P1 & Adaptive complexity ratio adjustment \\
SIM-CMP-004 & P2 & Handle missing historical data \\
SIM-CMP-005 & P2 & Respect complexity bounds \\
\bottomrule
\end{tabularx}

\subsection{Critic evaluation}

\begin{tabularx}{\linewidth}{@{}l l X@{}}
\toprule
\textbf{ID} & \textbf{Priority} & \textbf{Description} \\ \midrule
SIM-CRT-001 & P1 & Reject syntactically invalid SQL \\
SIM-CRT-002 & P1 & Reject semantically incorrect queries \\
SIM-CRT-003 & P1 & Accept valid query-answer pairs \\
SIM-CRT-004 & P1 & Adaptive threshold adjustment \\
SIM-CRT-005 & P2 & Handle critic model unavailability \\
\bottomrule
\end{tabularx}

\subsection{Diversity metrics}

\begin{tabularx}{\linewidth}{@{}l l X@{}}
\toprule
\textbf{ID} & \textbf{Priority} & \textbf{Description} \\ \midrule
SIM-DIV-001 & P1 & Calculate global diversity score \\
SIM-DIV-002 & P1 & Calculate local diversity score \\
SIM-DIV-003 & P2 & Coverage gap detection \\
SIM-DIV-004 & P2 & Diversity optimizer weight update \\
SIM-DIV-005 & P3 & Embedding computation for similarity \\
\bottomrule
\end{tabularx}

\subsection{Adaptive features}

\begin{tabularx}{\linewidth}{@{}l l X@{}}
\toprule
\textbf{ID} & \textbf{Priority} & \textbf{Description} \\ \midrule
SIM-ADP-001 & P1 & Fallback on adaptive complexity failure \\
SIM-ADP-002 & P1 & Fallback on adaptive threshold failure \\
SIM-ADP-003 & P1 & Fallback on coverage gap failure \\
SIM-ADP-004 & P2 & Alert on high failure rate \\
SIM-ADP-005 & P2 & Log failures for investigation \\
\bottomrule
\end{tabularx}

\subsection{HANA type handling}

\begin{tabularx}{\linewidth}{@{}l l X@{}}
\toprule
\textbf{ID} & \textbf{Priority} & \textbf{Description} \\ \midrule
SIM-HANA-001 & P1 & Map HANA types to standard SQL \\
SIM-HANA-002 & P2 & Handle spatial types (\texttt{ST\_GEOMETRY}, \texttt{ST\_POINT}) \\
SIM-HANA-003 & P2 & Handle full-text types (TEXT) \\
SIM-HANA-004 & P3 & Handle unknown types gracefully \\
\bottomrule
\end{tabularx}

\subsubsection{SIM-HANA-002 implementation}

Spatial columns (\texttt{ST\_GEOMETRY}, \texttt{ST\_POINT}) are \emph{not}
rendered directly in generated SQL; the extractor maps them to a stable set
of SAP HANA spatial \emph{methods} that return scalar or text-serialisable
values. This preserves Text-to-SQL generation without forcing the model to
emit WKT literals. The mapping is:

\begin{table}[H]
\centering\footnotesize
\caption{HANA spatial type $\to$ scalar method mapping.}
\label{tab:sim-hana-spatial-mapping}
\begin{tabularx}{\linewidth}{@{}l l X@{}}
\toprule
\textbf{Source type} & \textbf{Rendered expression} & \textbf{Return type} \\ \midrule
\texttt{ST\_POINT} & \texttt{col.ST\_X(), col.ST\_Y()} & \textsc{double} \\
\texttt{ST\_GEOMETRY} & \texttt{col.ST\_AsWKT()} & \textsc{nvarchar} \\
\texttt{ST\_GEOMETRY} (filter) & \texttt{col.ST\_Intersects(:bbox)=1} & \textsc{tinyint} \\
\texttt{ST\_GEOMETRY} (distance) & \texttt{col.ST\_Distance(:point, 'meter')} & \textsc{double} \\
Unknown spatial subtype & \texttt{col.ST\_AsWKT()} (fallback) & \textsc{nvarchar} \\
\bottomrule
\end{tabularx}
\end{table}

\noindent
Acceptance criteria for SIM-HANA-002:

\begin{enumerate}
  \item Given a table with an \texttt{ST\_POINT} column, the taxonomy
        builder registers \emph{two} synthetic projection factors
        (\texttt{.ST\_X()}, \texttt{.ST\_Y()}) and \emph{one} predicate
        factor (\texttt{.ST\_Distance(:point, 'meter') < :r}).
  \item Generated SQL passes \texttt{EXPLAIN PLAN} against the HANA
        test harness in \path{tests/fixtures/simula/hana-spatial/}
        (four fixture tables: \texttt{STORES}, \texttt{DELIVERY\_ZONES},
        \texttt{SENSORS}, \texttt{ROUTES}).
  \item When the column SRID is unknown, the extractor's behaviour is
        governed by the \texttt{spatial\_unknown\_srid\_policy}
        configuration flag (default \texttt{warn}; set in
        \path{configs/simula/schema-extraction.yaml}):
        \begin{itemize}
          \item \texttt{warn} (default): emit \texttt{ST\_AsWKT()},
                write a warning to
                \path{logs/simula/schema-extraction.log}, still
                generate the example, and flag
                \texttt{critic\_metadata.spatial\_srid\_unknown=true}.
                Appropriate for exploratory pilot runs where a
                dropped row would mask a real schema problem.
          \item \texttt{strict}: abort schema extraction with a
                non-zero exit code and an actionable error pointing
                at the offending column. Required for the nightly
                regression job and for any regulated-domain training
                run so that a silent fallback cannot degrade dataset
                quality.
          \item \texttt{skip}: drop the column from the taxonomy
                (the extractor logs which columns were skipped).
                Used when the column is known to be non-essential
                (e.g.\ audit-only spatial fields).
        \end{itemize}
        The policy is recorded in the reproducibility manifest
        (\texttt{manifest.json}, field
        \texttt{spatial\_unknown\_srid\_policy}) so the dataset
        provenance is unambiguous.
  \item The P2 test case passes if a 50-example generation run over
        the spatial fixture produces at least 40 syntactically valid
        queries (80\,\%) and zero queries containing raw WKT literals.
  \item Under \texttt{strict} policy, running the extractor against
        the \texttt{STORES\_BAD\_SRID} fixture must produce a
        non-zero exit code and write no taxonomy output --- this is
        asserted by \texttt{test\_spatial\_strict\_mode}.
\end{enumerate}

\noindent
Reference implementation:
\path{src/training/pipeline/hana_type_mapper.py::map_spatial_column()}.
Policy enforcement: \path{src/training/pipeline/schema_extractor.py::_apply_spatial_policy()}.

\section{Integration test suite}

\begin{lstlisting}[basicstyle=\ttfamily\small]
pytest tests/integration/simula/ -v --tb=short
\end{lstlisting}

Expected duration: 30 minutes.

\subsection{End-to-end scenarios}

\begin{enumerate}
  \item \textbf{Full generation run}: Generate 100 examples from taxonomy
  \item \textbf{Best-of-N selection}: Verify top candidate selection
  \item \textbf{Critic filtering}: Verify rejection of invalid examples
  \item \textbf{Diversity optimization}: Verify coverage improvement
  \item \textbf{Reproducibility}: Same seed produces same outputs
\end{enumerate}

\section{Performance benchmarks}

All p95 targets in Table~\ref{tab:sim-perf-targets} are measured on the
\emph{reference hardware profile} defined in
Section~\ref{sec:sim-perf-hardware}. Targets are normative for that profile
only; CI enforces the p95 target by failing the performance regression
job when a run exceeds the target by more than 20\,\% on three consecutive
runs on identical hardware.

\begin{table}[htbp]
\centering
\caption{Simula performance targets on reference hardware.}
\label{tab:sim-perf-targets}
\footnotesize
\begin{tabularx}{\linewidth}{@{}l r r X@{}}
\toprule
\textbf{Operation} & \textbf{p95 Target} & \textbf{Throughput} & \textbf{Notes} \\ \midrule
Example generation & 10,000 ms & 500/hr & Including LLM call \\
Critic evaluation & 2,000 ms & 1,000/hr & Per candidate \\
Diversity calculation & 100 ms & 10,000/hr & Per batch \\
Full pipeline (100 ex) & 30 min & --- & End-to-end \\
\bottomrule
\end{tabularx}
\end{table}

\subsection{Reference hardware profile}
\label{sec:sim-perf-hardware}

\begin{description}
  \item[GPU.] 1$\times$ NVIDIA A100 40\,GB SXM4 \emph{or} L40S 48\,GB
        (either satisfies the p95 targets in
        Table~\ref{tab:sim-perf-targets}). Development laptops with
        consumer GPUs (e.g.\ T4 15\,GB) are \emph{not} in scope.
  \item[Model.] \texttt{Qwen2.5-32B-Instruct} served via vLLM 0.6.x,
        tensor-parallel degree 1, \texttt{max\_model\_len=8192},
        \texttt{gpu\_memory\_utilization=0.90}.
  \item[Generation concurrency.] Batch size 1 per LLM call for the
        single-operation p95 target; pipeline throughput
        (500~examples/hour) measured at \texttt{concurrency=4} with
        four async workers sharing the vLLM server.
  \item[CPU / RAM.] 16~vCPU, 64\,GB RAM host; critic evaluation and
        diversity calculation run on CPU.
  \item[Storage.] NVMe SSD, $>$500\,MB/s sequential read; taxonomy
        loads cached in memory.
  \item[Network.] Co-located LLM server; round-trip latency
        $<$0.5\,ms (LLM on the same host).
\end{description}

\noindent
When the pipeline is run against a \emph{remote} LLM endpoint (e.g.\ a
shared model gateway), add the measured one-way network latency to each
\texttt{Example generation} p95 --- the LLM-call component is
$\sim$9.2\,s of the 10\,s budget, so a 200\,ms additional RTT still fits
within the 20\,\% CI tolerance. Hardware profile updates must be
approved by the Simula engineering lead and re-baselined against all
targets in Table~\ref{tab:sim-perf-targets}.

\section{Quality metrics tests}

\begin{tabularx}{\linewidth}{@{}l X@{}}
\toprule
\textbf{Test} & \textbf{Description} \\ \midrule
SIM-QM-001 & Rejection rate within expected range (5--15\%) \\
SIM-QM-002 & Diversity score meets target ($\geq$0.5) \\
SIM-QM-003 & Coverage at level 2 meets target ($\geq$0.8) \\
SIM-QM-004 & Complexity distribution matches target ratio \\
\bottomrule
\end{tabularx}

\section{CI/CD integration}

\begin{lstlisting}[basicstyle=\ttfamily\small,caption={GitHub Actions workflow excerpt.}]
simula-tests:
  runs-on: ubuntu-latest
  steps:
    - uses: actions/checkout@v4
    - name: Run Simula tests
      run: |
        pytest tests/unit/simula/ -v --cov=src/simula
        pytest tests/integration/simula/ -v --timeout=3600
    - name: Check coverage
      run: |
        coverage report --fail-under=85
\end{lstlisting}

\section{Acceptance criteria}

Tests pass when:
\begin{enumerate}
  \item All P1 test cases pass (100\%)
  \item P2 test cases pass ($\geq$95\%)
  \item P3 test cases pass ($\geq$90\%)
  \item Line coverage $\geq$85\%
  \item All fallback paths exercised
  \item Quality metrics within expected ranges
\end{enumerate}